\documentclass[10pt,a4paper]{amsart}
\usepackage[margin=19mm]{geometry}
\usepackage{amsmath,amssymb,mathtools}
\usepackage[scr=rsfs,cal=euler]{mathalfa}
\usepackage{xfrac}
\usepackage[unicode,hidelinks,bookmarks=false]{hyperref}
\newcommand{\ie}{i.e.\ }
\newcommand{\cf}{cf.\ }
\newcommand{\resp}{respectively\ }
\newcommand{\Subsection}{\subsection}
\providecommand{\nobreakdash}{\nobreak\hbox{-}}
\setlength{\emergencystretch}{2em}
\multlinegap0pt
\usepackage{bm}
\usepackage{bbm}
\usepackage{dsfont}
\usepackage{xspace}
\usepackage{cancel}
\usepackage{mathtools}
\usepackage{amsthm}
\makeatletter
\@ifundefined{lemma}{\newtheorem{lemma}{Lemma}}{}
\@ifundefined{prop}{\newtheorem{prop}{Proposition}}{}
\makeatother
\newcommand{\LC}{\left(}
\newcommand{\RC}{\right)}
\newcommand{\R}{\mathbb{R}}
\newcommand{\floor}[1]{\lfloor #1 \rfloor}
\title[The higher-order fractional Schr\"odinger equation with nonl]{The higher-order fractional Schr\"odinger equation with nonlinear local perturbations: Uniqueness}
\author{Giovanni Covi, Ru-Yu Lai, Lili Yan}
\date{Source: arXiv:2511.05384v1 (CC BY 4.0)}
\begin{document}
\maketitle

\noindent\emph{Source and licence.} The higher-order fractional Schr\"odinger equation with nonlinear local perturbations: Uniqueness. Version arXiv:2511.05384v1 (CC BY 4.0), distributed under the Creative Commons Attribution 4.0 International licence, as recorded in the arXiv licence metadata for this exact version and shown on the abstract page https://arxiv.org/abs/2511.05384v1. Full rights notice: licenses/arxiv-2511.05384-CC-BY-4.0.txt.\par\smallskip

\noindent\textit{Editorial scope.} Counted unit: the Lemma whose source label is \texttt{prop\_est\_U\&Tilde T} in the article (source0.tex lines 1093--1099, source label \texttt{prop\_est\_U\&Tilde T}), delivered together with the article's own complete original proof of it (source0.tex lines 1100--1182, one \texttt{proof} environment of 83 lines). No mathematics is written, repaired or re-derived: statement and proof are the article's own text line for line. Required context retained uncounted: eq\_linear (lines 229--234); prop\_Holder\_linear (lines 272--283); EQU\_U general (lines 1001--1006); EST:U\_j (lines 1008--1015); eq\_T\_k (lines 1038--1045); eq\_est\_V\&T (lines 1065--1069). Every definition, display and label of those lines is retained unchanged. Omitted from this package and not redistributed: 1--228, 235--271, 284--1000, 1007--1007, 1016--1037, 1046--1064, 1070--1092, 1183--1199. Nothing inside a retained excerpt is omitted or altered: every retained body is delivered exactly as the article's lines are written, so no omission receipt of any kind accompanies this package. Citations: the retained excerpts carry no citation command at all, so no bibliography entry of the article is transcribed and no reference is redistributed. Reference adaptation: none was needed and none was made; every reference inside the delivered unit resolves inside it, so no unresolved reference or citation is produced. Printed slips kept literal: no printed word, symbol, hypothesis or number of the counted statement or of its proof was altered. Layout and class: the article's own class is \texttt{amsart} with packages including amsfonts, amssymb, mathtools, hyperref, comment, color, xcolor, a4, geometry, todonotes, appendix, mdframed, url; the wrapper is the lane's pinned \texttt{amsart} editorial wrapper at 10pt on A4, which restates the theorem-like environments the retained text needs and adds the article's own macro definitions that the retained text needs (\texttt{\textbackslash LC}, \texttt{\textbackslash R}, \texttt{\textbackslash RC}, \texttt{\textbackslash floor}), with extra wrapper packages (bm, bbm, dsfont, xspace, cancel, mathtools). The delivered numbering is the unit's own. Assets: no retained span contains a figure environment, a graphics-inclusion command, a PDF-page inclusion, a table or a TikZ picture; no image file is copied into this package, the package declares no asset dependency and no image file is redistributed with it. Licence: version arXiv:2511.05384v1 of this article is distributed under the Creative Commons Attribution 4.0 International licence, as recorded in the arXiv licence metadata for this exact version and shown on the abstract page https://arxiv.org/abs/2511.05384v1. A full rights notice is delivered at licenses/arxiv-2511.05384-CC-BY-4.0.txt. Characters: every retained excerpt is pure ASCII, so the delivered engine, XeLaTeX with its default Latin Modern text font, reports no missing character.
\begin{equation}\label{eq_linear}
\begin{cases}
(-\Delta)^s v + q(x)v = F &\hbox{ in } \Omega,\\
v = f &\hbox{ in } \Omega_e,
\end{cases}
\end{equation}

\begin{prop}\label{prop_Holder_linear}
Let $\Omega\subset\R^n$, $n\geq 2$, be a bounded and smooth open domain. Let $s\in \R^+\setminus \mathbb{Z}$, and assume $\floor{s}>n/2$. Suppose that $f\in C_c^\infty(\Omega_e)$, $F\in L^\infty(\Omega)$, and  
$q\in L^\infty(\Omega)$ such that $q\ge 0$ in $\Omega$.
Let $v\in H^s(\R^n)$ solve \eqref{eq_linear}.
Then $v\in C^{s}(\R^n)$ and it satisfies
\[
\|v\|_{C^{s}(\R^n)}
\le
C \LC  \|F\|_{L^\infty(\Omega)}+ \|f\|_{C_c^\infty(\Omega_e)} \RC,
\]
for a constant $C>0$ independent of $v, F$ and $f$. 
\end{prop}

\begin{align}\label{EQU_U general}
\begin{cases}
[(-\Delta)^s +q(x) ]U_j  = -\widetilde{\mathcal{T}}_j &\hbox{ in } \Omega,\\
U_j = 0 &\hbox{ in } \Omega_e,
\end{cases}
\end{align}

\begin{equation}\label{EST:U_j}  
\begin{aligned}
&\|U_1\|_{C^s(\R^n)}+\|U_1\|_{H^s(\R^n)}
\le  C\|\mathbf{P}(u_\varepsilon) \|_{L^\infty(\Omega)},\\
&\|U_j\|_{C^s(\R^n)}+\|U_j\|_{H^s(\R^n)}
\le  C\|\widetilde{\mathcal{T}}_{j}\|_{L^\infty(\Omega)}. 
\end{aligned}
\end{equation}

\begin{equation}\label{eq_T_k}
\begin{aligned}
\mathcal T_{k} 
& =
\sum^{k-1}_{\ell=1} \sum_{\substack{ \gamma_1,\ldots, \gamma_{\ell+1}\in \mathbb{N} : \\ \sum_{j=1}^{\ell+1}\gamma_j = k }}  \bigg( \prod_{j=1}^\ell\sum_{|\beta_j|=\gamma_j}  \frac{\varepsilon^{\beta_j} w_{\beta_j }}{\beta_j!} \bigg)    P_\ell(x,D)\bigg( \sum_{|\beta_{\ell+1}|=\gamma_{\ell+1}}\frac{\varepsilon^{\beta_{\ell+1}}w_{\beta_{\ell+1}}}{\beta_{\ell+1}!}\bigg) \\
& = \sum^{k-1}_{\ell=1} \sum_{\substack{ \gamma_1,\ldots, \gamma_{\ell+1}\in \mathbb{N} : \\ \sum_{j=1}^{\ell+1}\gamma_j = k }}  \bigg( \prod_{j=1}^\ell V_{\gamma_j} \bigg)    P_\ell(x,D)V_{\gamma_{\ell+1}}.
\end{aligned}
\end{equation}

\begin{equation}\label{eq_est_V&T}
\|\mathcal{T}_k\|_{L^\infty(\Omega)} = \mathcal{O}(\|\varepsilon\cdot f\|^k_{C_c^\infty(\Omega_e)}),\quad
\|V_k\|_{C^s(\R^n)}+\|V_k\|_{H^s(\R^n)}
= \mathcal{O}(\|\varepsilon\cdot f\|^k_{C_c^\infty(\Omega_e)}).
\end{equation}

\begin{lemma}\label{prop_est_U&Tilde T}
For $0\le k\le K$, we have
\begin{equation}\label{eq_est_U&Tilde T}
\|\widetilde{\mathcal{T}}_k\|_{L^\infty(\Omega)} = \mathcal{O}(\|\varepsilon\cdot f\|^{k+1}_{C_c^\infty(\Omega_e)}),\quad
\|U_k\|_{C^s(\R^n)}+ \|U_k\|_{H^s(\R^n)}= \mathcal{O}(\|\varepsilon\cdot f\|^{k+1}_{C_c^\infty(\Omega_e)}).
\end{equation}
\end{lemma}

\begin{proof}
We proceed again by induction. From the definition $U_0=u_\varepsilon$, it is clear that $U_0$ is of order $1$.
As $\mathcal{T}_0=\mathcal{T}_1=0$, we have $\widetilde{\mathcal{T}}_0=\widetilde{\mathcal{T}}_1=\mathbf{P}(U_0)=\sum_{\ell = 1}^{K-1}U_0^{\ell}P_\ell(x,D)U_0$.
From the estimate
\begin{align*}
\|\mathbf{P}(U_0)\|_{L^\infty(\Omega)}\leq \sum_{\ell = 1}^{K-1}\|U_0\|_{C^{s}(\R^n)}^{\ell}\|P_\ell(x,D)U_0\|_{C(\R^n)} \le \sum_{\ell = 1}^{K-1}\|U_0\|_{C^{s}(\R^n)}^{\ell+1} = \mathcal{O}(\|\varepsilon\cdot f\|^{2}_{C_c^\infty(\Omega_e)}),
\end{align*}	
we deduce
$$
\widetilde{\mathcal{T}}_0=\widetilde{\mathcal{T}}_1 = \mathcal{O}(\|\varepsilon\cdot f\|^{2}_{C_c^\infty(\Omega_e)}).
$$
With this, since $U_1$ solves \eqref{EQU_U general}, Proposition~\ref{prop_Holder_linear} gives
\begin{align}
\|U_1\|_{C^s(\R^n)}+ \|U_1\|_{H^s(\R^n)}
\le  C\|\widetilde{\mathcal{T}}_1 \|_{L^\infty(\Omega)}=\mathcal{O}(\|\varepsilon \cdot f\|^{2}_{C^\infty_c(\Omega_e)}).
\end{align}	


By the induction argument, assume that for a positive integer $J\in \{2,\ldots, K-1\}$, the following holds for all $0\leq k\leq J$:
$$
\widetilde {\mathcal{T}}_{k} =  \mathcal{O}(\|\varepsilon\cdot f\|^{k+1}_{C^\infty_c(\Omega_e)}),\quad \hbox{ and }\quad\| U_{k}\|_{C^{s}(\R^n)}+ \|U_k\|_{H^s(\R^n)} = \mathcal{O}(\|\varepsilon\cdot f\|^{k+1}_{C^\infty_c(\Omega_e)}).
$$
We will prove that $\widetilde {\mathcal{T}}_{J+1}$ is of order $J+2$. Using \eqref{eq_T_k}, we derive
\begin{align*}
\widetilde{\mathcal{T}}_{J+1} & = \mathbf{P}(U_0) - \sum_{b=2}^{J+1}\mathcal T_b = \sum_{\ell=1}^{K-1}U_0^\ell P_\ell (x,D)U_0 - \sum_{b=1}^{J}\mathcal T_{b+1}
\\ & =
\sum_{\ell=1}^{K-1}U_0^\ell P_\ell(x,D) U_0 - \sum_{b=1}^{J}\sum^{b}_{\ell=1} \sum_{\substack{ \gamma_1,\ldots, \gamma_{\ell+1}\in \mathbb{N}  : \\ \sum_{j=1}^{\ell+1}\gamma_j = b+1 }}  \bigg( \prod_{j=1}^\ell V_{\gamma_j} \bigg) P_\ell(x,D) V_{\gamma_{\ell+1}} 
\\ & =
\sum_{\ell = J+1}^{K-1}U_0^{\ell}P_\ell(x,D)U_0 +  \sum^{J}_{\ell=1} \left(U_0^ {\ell}P_\ell(x,D)U_0 - \underbrace{\sum_{b=\ell}^{J}\sum_{\substack{ \gamma_1,\ldots, \gamma_{\ell+1}\in \mathbb{N} : \\ \sum_{j=1}^{\ell+1}\gamma_j = b+1 }}  \bigg( \prod_{j=1}^\ell V_{\gamma_j} \bigg)    P_\ell(x,D)V_{\gamma_{\ell+1}} }_{I_\ell}\right).
\end{align*}	
Since $\|U_0\|_{C^s(\R^n)} =\mathcal{O}(\|\varepsilon\cdot f\|^{1}_{C_c^\infty(\Omega_e)})$, it follows that the first term satisfies 
\[
\left\| \sum_{\ell = J+1}^{K-1}U_0^{\ell}P_\ell(x,D)U_0 \right\|_{C^s(\R^n)} =\mathcal{O}(\|\varepsilon\cdot f\|^{J+2}_{C_c^\infty(\Omega_e)}). 
\]
It remains to show that the second term in $\widetilde{\mathcal{T}}_{J+1}$ is also of order $J+2$.

Denote $\sigma_0 := 0$ and $\sigma_j := \sum_{s=1}^j\gamma_s$, $1\le j\le \ell+1$. Then $\gamma_j = \sigma_j - \sigma_{j-1}$, and also $\sigma_{l+1}=b+1$. Since $\gamma_j\ge 1$, we can split and reorder the summation 
\begin{align*}
\mathrm{I}_{\ell} &=\sum_{b=\ell}^J\sum_{\substack{ \gamma_1,\ldots, \gamma_{\ell+1}\in \mathbb{N}  : \\ \sum_{j=1}^{\ell+1}\gamma_j = b+1 }}\bigg( \prod_{j=1}^\ell V_{\sigma_j-\sigma_{j-1}} \bigg)    P_\ell(x,D)V_{{\sigma_{\ell+1}-\sigma_{\ell}}}  \\
& = 
\sum_{b=\ell}^J\sum_{\sigma_1 = 1}^{(b-\ell+1)} \ldots\sum_{\sigma_n = 1 + \sigma_{n-1}}^{(b-\ell+n)} \ldots \sum_{\sigma_{\ell}= 1+\sigma_{\ell-1} }^b\bigg( \prod_{j=1}^\ell V_{\sigma_j-\sigma_{j-1}} \bigg)    P_\ell(x,D)V_{{b+1-\sigma_{\ell}}} 
\\ & =
\sum_{\sigma_1 = 1}^{(J-\ell+1)} ... \sum_{\sigma_n = 1 + \sigma_{n-1}}^{(J-\ell+n)} \ldots \sum_{\sigma_\ell = 1 + \sigma_{\ell-1}}^{J} \sum_{b=\sigma_\ell}^J \bigg( \prod_{j=1}^\ell V_{\sigma_j-\sigma_{j-1}} \bigg)    P_\ell(x,D)V_{{b+1-\sigma_{\ell}}} \\
&=  \sum_{\sigma_1 = 1}^{(J-\ell+1)} \ldots \sum_{\sigma_n = 1 + \sigma_{n-1}}^{(J-\ell+n)}\ldots \sum_{\sigma_\ell = 1 + \sigma_{\ell-1}}^{J}  \bigg( \prod_{j=1}^\ell V_{\sigma_j-\sigma_{j-1}} \bigg)    P_\ell(x,D) \left(\sum_{b=\sigma_\ell}^JV_{b+1-\sigma_\ell}\right).
\end{align*}	
We observe that
\begin{align*}
\left(\sum_{b=\sigma_\ell}^JV_{b+1-\sigma_\ell}\right) = \sum_{b=\sigma_\ell}^J(U_{b-\sigma_\ell}- U_{b+1-\sigma_\ell}) =   U_{0} - U_{J+1-\sigma_\ell},
\end{align*}
and as $\sigma_0=0$, we have by \eqref{eq_est_V&T} that
\begin{align*}
\prod_{j=1}^\ell V_{\sigma_j-\sigma_{j-1}}  =  \prod_{j=1}^\ell \underbrace{(U_{\sigma_j-\sigma_{j-1}-1}-U_{\sigma_j-\sigma_{j-1}})}_{\mathcal{O}(\|\varepsilon\cdot f\|^{\sigma_j-\sigma_{j-1}}_{C_c^\infty(\Omega_e)})}= \mathcal{O}(\|\varepsilon\cdot f\|^{\sigma_\ell}_{C_c^\infty(\Omega_e)}).
\end{align*}	
It leads to 
\begin{align*}
\mathrm{I}_\ell 
&=  \sum_{\sigma_1 = 1}^{(J-\ell+1)} \ldots \sum_{\sigma_n = 1 + \sigma_{n-1}}^{(J-\ell+n)} \ldots \sum_{\sigma_\ell = 1 + \sigma_{\ell-1}}^{J}  \underbrace{\bigg( \prod_{j=1}^\ell V_{\sigma_j-\sigma_{j-1}} \bigg)}_{ \mathcal{O}(\|\varepsilon\cdot f\|^{\sigma_\ell}_{C_c^\infty(\Omega_e)})}    P_\ell(x,D) U_0\\
&\quad - \sum_{\sigma_1 = 1}^{(J-\ell+1)}\ldots \sum_{\sigma_n = 1 + \sigma_{n-1}}^{(J-\ell+n)}\ldots \sum_{\sigma_\ell = 1 + \sigma_{\ell-1}}^{J}  \underbrace{\bigg( \prod_{j=1}^\ell V_{\sigma_j-\sigma_{j-1}} \bigg)}_{ \mathcal{O}(\|\varepsilon\cdot f\|^{\sigma_\ell}_{C_c^\infty(\Omega_e)})}    P_\ell(x,D) \underbrace{U_{J+1-\sigma_\ell}}_{\mathcal{O}(\|\varepsilon\cdot f\|^{J+2-\sigma_\ell}_{C_c^\infty(\Omega_e)})},\\
&= \underbrace{\sum_{\sigma_1 = 1}^{(J-\ell+1)} \ldots \sum_{\sigma_n = 1 + \sigma_{n-1}}^{(J-\ell+n)} \ldots \sum_{\sigma_\ell = 1 + \sigma_{\ell-1}}^{J}  \underbrace{\bigg( \prod_{j=1}^\ell V_{\sigma_j-\sigma_{j-1}} \bigg)}_{ \mathcal{O}(\|\varepsilon\cdot f\|^{\sigma_\ell}_{C_c^\infty(\Omega_e)})}    P_\ell(x,D) U_0 }_{\mathrm{I}_{\ell,1}} +\mathcal{O} (\|\varepsilon\cdot f\|^{J+2}_{C_c^\infty(\Omega_e)}),
\end{align*}	
where the order of $U_{J+1-\sigma_\ell}$ is $J+2-\sigma_\ell$ by the induction hypothesis.


Now, we estimate 
\begin{align*}
\mathrm{I}_{\ell,1} 
&=  \sum_{\sigma_1 = 1}^{(J-\ell+1)}\ldots \sum_{\sigma_n = 1 + \sigma_{n-1}}^{(J-\ell+n)}\ldots \sum_{\sigma_{\ell-1} = 1 + \sigma_{\ell-2}}^{J-1}  \bigg( \prod_{j=1}^{\ell-1} V_{\sigma_j-\sigma_{j-1}} \bigg) \underbrace{\bigg(\sum_{\sigma_\ell= 1 +  \sigma_{\ell-1}}^{J} V_{\sigma_\ell-\sigma_{\ell-1}} \bigg)}_{U_0-U_{J-\sigma_{\ell-1}}}  P_\ell(x,D) U_{0}  \\
&=  \sum_{\sigma_1 = 1}^{(J-\ell+1)}\ldots \sum_{\sigma_n = 1 + \sigma_{n-1}}^{(J-\ell+n)}\ldots \sum_{\sigma_{\ell-1} = 1 + \sigma_{\ell-2}}^{J-1}  \bigg( \prod_{j=1}^{\ell-1} V_{\sigma_j-\sigma_{j-1}} \bigg)  U_0  P_\ell(x,D) U_{0} \\
&\quad - \sum_{\sigma_1 = 1}^{(J-\ell+1)}\ldots \sum_{\sigma_n = 1 + \sigma_{n-1}}^{(J-\ell+n)}\ldots \sum_{\sigma_{\ell-1} = 1 + \sigma_{\ell-2}}^{J-1}  \bigg( \prod_{j=1}^{\ell-1} V_{\sigma_j-\sigma_{j-1}} \bigg)   U_{J-\sigma_{\ell-1}}   P_\ell(x,D) U_{0},
\end{align*}	
where the same argument yields that the second term of $I_{\ell,1}$ above is also of order $J+2$.
We continue this procedure, which eventually leads us to 
\begin{align*}
\mathrm{I}_{\ell} = U_0^\ell P_{\ell}(x,D) U_0 + \mathcal{O}(\|\varepsilon\cdot f\|^{J+2}_{C_c^\infty(\Omega_e)}).
\end{align*}
Therefore, we obtain 
\begin{align*}
\widetilde{\mathcal{T}}_{J+1} = \mathcal{O}(\|\varepsilon\cdot f\|^{J+2}_{C_c^\infty(\Omega_e)})+\sum^{J}_{\ell=1} \left(U_0^ {\ell}P_\ell(x,D)U_0 - \mathrm{I}_{\ell}\right) =  \mathcal{O}(\|\varepsilon\cdot f\|^{J+2}_{C_c^\infty(\Omega_e)}),
\end{align*}
and also \eqref{EST:U_j} implies 
$U_{J+1}=  \mathcal{O}(\|\varepsilon\cdot f\|^{J+2}_{C_c^\infty(\Omega_e)})$. This completes the proof.

\end{proof}

\end{document}
